% Part: methods
% Chapter: induction
% Section: induction-on-N

\documentclass[../../../include/open-logic-section]{subfiles}

\begin{document}

\olfileid{mth}{ind}{inN}

\olsection{~$\Nat$ वरील विगमन}

सर्वात साध्या रूपात विगमन ही सर्व नैसर्गिक संख्यांसाठी निष्कर्ष सिद्ध करण्याची पद्धती आहे. $0$ पासून सुरुवात करून वारंवार~$1$ वाढवत गेल्यास प्रत्येक नैसर्गिक संख्येपर्यंत अखेर पोहोचता येते, या तथ्याचा ती वापर करते. म्हणून एखादी गोष्ट प्रत्येक संख्येसाठी खरी आहे हे सिद्ध करण्यासाठी (1) ती $0$ साठी खरी आहे हे प्रस्थापित करता येते आणि (2) ती एखाद्या~$n$ संख्येसाठी खरी असेल तेव्हा पुढच्या~$n+1$ संख्येसाठीही खरी असते हे दाखवता येते. “~$n$ संख्येत $P$ हा गुणधर्म आहे” याचा संक्षेप $P(n)$ असा करू; तसेच “~$k$ संख्येत $P$ हा गुणधर्म आहे” याचा $P(k)$, इत्यादी. मग सर्व $n \in \Nat$ साठी $P(n)$ सिद्ध करणाऱ्या विगमनाधारित सिद्धतेत पुढील गोष्टी असतात:
\begin{enumerate}
\item $P(0)$ ची सिद्धता; आणि
\item कोणत्याही~$k$ साठी, $P(k)$ असेल तर $P(k+1)$ असते याची सिद्धता.
\end{enumerate}
हे पूर्ण स्पष्ट करण्यासाठी (1) आणि~(2) दोन्ही आहेत असे समजा. मग (1) वरून $P(0)$ सत्य आहे असे समजते. ~(2) देखील असल्यामुळे विशेषतः $P(0)$ असेल तर $P(0+1)$, म्हणजे $P(1)$, असते हे समजते. “कोणत्याही~$k$ साठी, $P(k)$ असेल तर $P(k+1)$” या सामान्य विधानात~$k$ च्या जागी~$0$ ठेवल्यावर हे मिळते. म्हणून विध्यात्मक अनुमानाने~$P(1)$ मिळते. पुन्हा (2) मध्ये, आता~$k$ च्या जागी $1$ घेतल्यावर, $P(1)$ असेल तर~$P(2)$ मिळते. ~$P(1)$ नुकतेच प्रस्थापित केलेले असल्यामुळे विध्यात्मक अनुमानाने~$P(2)$ मिळते. असेच पुढे चालू राहते. कोणत्याही~$n$ संख्येसाठी ही क्रिया $n$~वेळा केल्यावर अखेर~$P(n)$ मिळते. म्हणून (1) आणि~(2) मिळून कोणत्याही $n \in \Nat$ साठी~$P(n)$ प्रस्थापित करतात.

एक उदाहरण पाहू. $n$ फासे टाकल्यावर किती वेगवेगळ्या बेरजा मिळू शकतात हे शोधायचे आहे असे समजा. थोडे विचित्र वाटले, तरी $0$ फाशांपासून सुरुवात करू. फासेच नसतील, तर “टाकून” मिळणारी एकच बेरीज आहे: कोणतेही ठिपके नाहीत; म्हणजे बेरीज~$0$. म्हणून बेरजेच्या वेगवेगळ्या शक्यतांची संख्या~$1$ आहे. एकच फासा असेल, म्हणजे $n=1$, तर $1$ ते~$6$ अशी सहा मूल्ये शक्य आहेत. दोन फाशांनी $2$ ते $12$ पर्यंत कोणतीही बेरीज मिळू शकते; म्हणजे $11$ शक्यता. तीन फाशांनी $3$ ते $18$ पर्यंत कोणतीही संख्या मिळू शकते; म्हणजे $16$ वेगवेगळ्या शक्यता. $1$, $6$, $11$, $16$: एक नमुना दिसतो. कदाचित उत्तर $5n+1$ असेल? अर्थात, $5n+1$ ही कमाल शक्य संख्या आहे. कारण $n$ फाशांनी मिळणाऱ्या सर्वात लहान $n$ मूल्यापासून---सर्व फाशांवर $1$---सर्वात मोठ्या $6n$ मूल्यापर्यंत---सर्व फाशांवर $6$---फक्त $5n+1$ संख्या आहेत.

\begin{thm}
$n$ फाशांनी $n$ ते $6n$ यांमधील सर्व $5n+1$ शक्य मूल्ये मिळवता येतात.
\end{thm}

\begin{proof}
$P(n)$ हे विधान असू द्या: “$n$~फासे वापरून $n$ ते~$6n$ यांमधील कोणतीही संख्या मिळवता येते.” शून्य फाशांचे प्रकरण वर पाहिले आहे: रिक्त बेरीज शून्य असून एकच शक्यता असते. धन संख्यांसाठी विगमन वापरण्यास पुढील गोष्टी सिद्ध करू:
\begin{enumerate}
\item \emph{विगमनाचा आधार} $P(1)$: एकच फासा वापरून $1$ ते $6$ यांमधील कोणतीही संख्या मिळवता येते.
\item \emph{विगमनाची पायरी}: प्रत्येक धन नैसर्गिक संख्या $k$ साठी, $P(k)$ असेल तर~$P(k+1)$.
\end{enumerate}

(1) हे $6$ पृष्ठे असलेला फासा पाहून सिद्ध होते. त्याला सहाही पृष्ठे आहेत आणि $1$ ते~$6$ मधील प्रत्येक संख्या एका पृष्ठावर दिसते. म्हणून एक फासा वापरून $1$ ते~$6$ मधील कोणतीही संख्या मिळवता येते.

(2) सिद्ध करण्यासाठी सशर्त विधानाचे पूर्वांग, म्हणजे $P(k)$, गृहीत धरतो. या गृहीतकाला \emph{विगमन गृहीतक} म्हणतात. त्याचा वापर करून~$P(k+1)$ सिद्ध करतो. अवघड भाग असा: $k+1$ फाशांनी मिळणाऱ्या शक्य बेरजांचा विचार $k$~फाशांनी मिळणाऱ्या बेरजा आणि अतिरिक्त $k+1$ व्या फाशाचे मूल्य यांच्या आधारे करण्याचा मार्ग शोधायचा. विगमन गृहीतक वापरायचे असेल, तर असेच करावे लागते.

विगमन गृहीतकानुसार $k$~फाशांनी $k$ ते~$6k$ मधील कोणतीही संख्या मिळवता येते. ~$(k+1)$ व्या फाशावर~$1$ आला, तर एकूण बेरजेत $1$ वाढतो. म्हणून $k$~फासे टाकून नंतर $(k+1)$ व्या फाशावर~$1$ आणल्यास $k+1$ ते $6k+1$ मधील कोणतेही मूल्य मिळवता येते. काय उरले? $6k+2$ ते $6k+6$ ही मूल्ये. $k$ फाशांवर $6$ आणून आणि $(k+1)$ व्या फाशावर $2$ ते $6$ मधील संख्या आणून ही मूल्ये मिळवता येतात. या दोन्ही गोष्टी मिळून असे दाखवतात की $k+1$ फाशांनी $k+1$ ते $6(k+1)$ मधील कोणतीही संख्या मिळवता येते. म्हणजे~$P(k)$ हे विगमन गृहीतक वापरून~$P(k+1)$ सिद्ध केले आहे.
\end{proof}

अनेकदा वस्तूंच्या अनुक्रमाविषयी---संख्या, संच इत्यादी---काही सिद्ध करायचे असते आणि तो अनुक्रमच “विगमनाधारित” व्याख्येने दिलेला असतो. म्हणजे $(n+1)$ वी वस्तू $n$ व्या वस्तूच्या आधारे परिभाषित केलेली असते. अशा वेळी विगमन वापरतो. उदाहरणार्थ,~$n$ पर्यंतच्या नैसर्गिक संख्यांची बेरीज~$s_n$ अशी परिभाषित करता येते:
\begin{align*}
  s_0 & = 0\\
  s_{n+1} & = s_n + (n+1)
\end{align*}
या व्याख्येवरून असे मिळते:
\begin{align*}
  s_0 & = 0,\\
  s_1 & = s_0 + 1 && = 1,\\
  s_2 & = s_1 + 2 && = 1 + 2 = 3\\
  s_3 & = s_2 + 3 && = 1 + 2 + 3 = 6, \text{ इत्यादी.}
\end{align*}
आता विगमनाने $s_n = n(n+1)/2$ सिद्ध करता येते.

\begin{prop}
  $s_n = n(n+1)/2$.
\end{prop}

\begin{proof}
(1) $s_0 = 0\cdot(0 + 1)/2$ आणि (2) $s_k = k(k+1)/2$ असेल तर $s_{k+1} = (k+1)(k+2)/2$, हे सिद्ध करायचे आहे. (1) स्पष्ट आहे. (2) सिद्ध करण्यासाठी $s_k = k(k+1)/2$ हे विगमन गृहीतक घेऊ. त्याचा वापर करून $s_{k+1} = (k+1)(k+2)/2$ दाखवायचे आहे.

$s_{k+1}$ काय आहे? व्याख्येनुसार $s_{k+1} = s_k + (k+1)$. विगमन गृहीतकानुसार $s_k = k(k+1)/2$. आधीच्या समीकरणात ही किंमत ठेवता येते; मग अपूर्णांकांचे थोडे अंकगणित पुरते:
\begin{align*}
    s_{k+1} & = \frac{k(k+1)}{2} + (k+1) = {}\\
    & = \frac{k(k+1)}{2} + \frac{2(k+1)}{2} = {}\\
    & = \frac{k(k+1) + 2(k+1)}{2} = {}\\
    & = \frac{(k+2)(k+1)}{2}.
\end{align*}
\end{proof}

महत्त्वाचा धडा असा: विगमनाधारित व्याख्येने दिलेल्या $a_n$ अनुक्रमाविषयी काही सिद्ध करत असाल, तर विगमन हा स्वाभाविक मार्ग आहे. तसे नसले तरी---फाशांच्या शक्य बेरजांच्या उदाहरणाप्रमाणे---~$k+1$ चे प्रकरण~$k$ च्या प्रकरणाशी काही रीतीने जोडता आले, तर विगमन वापरता येते.

\end{document}
